%قبل از شروع مسئله ترجیحاً چیزی ننویسید. اما اگر خواستید بنویسید، در قالب لیستی (فقرات و شمارش) بنویسید.
\شروع‌مسئله[12]  
% صورت اصلی مسئله را در این‌جا بنویسید.
\\
شنگدباو برای اینکه ناچار به اعتراف نادانی خود نشود از شما تعداد زیادی سؤال پرسیده و از شما می‌خواهد سریعتر جواب دهید. او دست پیش‌ را گرفته‌است تا پس نیفتد! بنابراین به پرسش‌هایش جواب دهید.
%اگر مسئله چند قسمت داشت به این صورت عمل کنید:
\شروع{زیرمسائل}
\زیرمسئله[3]
به ازای رشته
$baabaabasse$
آرایه پسوندی این رشته را بنویسید. (صرفاً پسوندها را به ترتیب لغت‌نامه‌ای ارائه دهید)

\زیرمسئله[3] 
برای همان رشته بخش قبل درخت پسوندی را بکشید (یعنی همه پسوند‌ها را در یک 
\lr{tri}
اضافه کنید و درخت را بنویسید.
)
%\زیرمسئله[0] % مسئلۀ ۱پ
%\زیرمسئله[0] % مسئلۀ ۱ت
\پایان{زیرمسائل}

حالا که به سؤال‌های بدیهی شنگدباو پاسخ داده‌اید نوبت شماست! 
\lr{compress}
کردن یک درخت پسوندی را به صورت زیر تعریف می‌کنیم:
\begin{itemize}
	\item 
	برای 
	\lr{compress}
	کردن یک درخت پسوندی اینگونه عمل می‌کنیم: فرض کنید یال‌ها به جای اینکه رویشان کاراکتر نوشته‌شده باشد یک رشته نوشته شده است. برای مثال اگر با طی مسیری از ریشه تا یک برگ به ترتیب رشته‌های
	$s_1, s_2, ..., s_k$
	را ببینیم این به این معناست که رشته‌
	$s_1 + s_2 + ... + s_k$
	یک پسوند از رشته ما است.
	
	در چنین درختی اگر رأسی تنها یک فرزند داشت و ریشه نبود می‌توانیم یال این رأس به فرزند و این رأس به پدرش را با هم ادغام کنیم تا یک یال بین پدرش و فرزندش به‌دست بیاید و رشته نوشته شده روی این یال برابر رشته حاصل از چسباندن یال بالایی به پایینی است.
	
	اگر این روند را ادامه دهیم تا هیچ رأس این‌چنینی‌‌ای دیگر موجود نباشد یعنی درخت را 
	\lr{compress}
	کرده‌ایم. (برای اینکه ذخیره‌سازی خود رشته حافظه اضافی مصرف نکند معمولا آن‌را با بازه‌ای نمایش می‌دهند 
	$[l,r]$
	که به معنای زیر رشته رشته‌ اصلی با شروع از $l$ و پایان $r$ است.
	)
\end{itemize}

سؤالات زیر را شما از شنگدباو می‌پرسید اما بعداً با خودتان فکر می‌کنید که آیا خودتان جواب چنین سؤالاتی را می‌دانید؟

برای اینکه خودتان را به خودتان ثابت کنید سؤالات زیر را حل کنید:

\شروع{زیرمسائل}
\زیرمسئله[4]
ثابت کنید هر رشته به طول $n$ درخت پسوندی 
\lr{compress}
شده‌اش حاوی 
$O(n)$
تا رأس و یال است.

\زیرمسئله[2]
در بخش قبل ثابت کردیم این درخت اندازه‌اش
$O(n)$
است اما در مورد پیدا کردن این درخت با همین مرتبه زمانی حرفی نزدیم.
یکی از الگوریتم‌های زیبای رشته‌ها به‌دست آوردن درخت پسوندی با 
$O(n)$
است! فرض کنید چنین الگوریتمی را بلدیم و با فرض اینکه تعداد حروف الفبا عددی ثابت است الگوریتمی از 
$O(n)$
برای پیدا کردن آرایه پسوندی از روی این درخت پسوندی ارائه دهید. 
%\زیرمسئله[0] % مسئلۀ ۱پ
%\زیرمسئله[0] % مسئلۀ ۱ت
\پایان{زیرمسائل}


\solution{

\begin{itemize}
\item [الف)]
$$aabaabasse$$
$$aabasse$$
$$abaabasse$$
$$abasse$$
$$asse$$
$$baabaabasse$$
$$baabasse$$
$$basse$$
$$e$$
$$se$$
$$sse$$

\item [ب)]
در شکل زیر به صورت 
\lr{compress}
شده درخت ذخیره شده‌است اما جواب‌هاییکه به صورت کامل هستند هم درست در نظر گرفته می‌شود.

\begin{center}
	\includegraphics [scale=0.5] {src/suffix_tree}
\end{center}

\item [پ)]
درخت پسوندی 
فشرده هر رأسی که رشته متناظرش یک پسوند نیست لااقل دو تا بچه‌ دارد زیرا در غیر اینصورت این رأس حذف می‌شود و یال فرزندش با یال پدرش ادغام می‌شود.
فرض کنید رئوسی که رشته متناظرش یک پسوند است را سیاه می‌کنیم و بقیه رئوس سفیداند. رئوس سفید لااقل دو تا فرزند دارند.

به استقرا ثابت می‌کنیم زیر درخت هر رأس که حال $B$ تا رأس سیاه است حداکثر اندازه‌اش $2B-1$ است. برای اثبات چنین چیزی فرض کنید ریشه سفید است، در اینصورت لا اقل دو تا فرزند دارد و فرض کنید فرزندان آن 
$u_1, ..., u_k$
است ($k \ge 1$) در اینصورت طبق فرض اسقرا داریم:
$$sz(v) = 1 + \sum_{1 \le i \le k} sz(u_i) \le 1 + \sum_{1 \le i \le k} (2B(u_i) - 1) = 2B(v) - k + 1 \le 2B(v) - 1$$
که در آن $sz(v)$ برابر اندازه زیر درخت $v$ است و $B(v)$ تعداد رئوس سیاه در زیر درخت $v$ است.

حالا اگر رأس مورد نظر سیاه باشد دیگر محدودیت $k>1$ را نداریم:
$$sz(v) = 1 + \sum_{1 \le i \le k} sz(u_i) \le 1 + \sum_{1 \le i \le k} (2B(u_i) - 1) = 2(B(v)-1) - k + 1  \le 2B(v) - 1$$
اما همانطور که می‌بینید همچنان فرض استقرا برقرار است.

 

\item [ت)]
کافیست روی درخت پیمایش $dfs$ کنیم و ترتیب بچه‌های هر کس بر اساس رشته روی یال فرزندان به صورت مرتب‌شده باشد.  هر موقع به رأسی رسیدیم که متناظر یک پسوند علامت گذاری شده صرفا آن‌را خروجی می‌دهیم. یه سادگی می‌توان دید رشته‌ها به ترتیب الفبایی خروجی داده می‌شوند.
\end{itemize}

}



%%%%%%%%%%%%%%%%%%%%%%%%%%%%%%%%%%%%%%%%%%%%%%%%%%%%%
% برخی دستورات مفید برای تمرین‌ها. %
%%%%%%%%%%%%%%%%%%%%%%%%%%%%%%%%%%%%%%%%%%%%%%%%%%%%%

% فرمول‌های نمونه:
%$f(x) = \Theta(x)$
%$T(x, y) \leq \log(x) + 2^y + \binom{2n}{n}$
% {\tt code\_function}


% لیست موارد بدون شماره‌گذاری:
%\شروع{فقرات}
%    \فقره مورد اول 
%        \شروع{فقرات}
%            \فقره مورد اول داخلی 
%                \شروع{فقرات}
%                    \فقره اولین مورد خیلی داخلی 
%                    \فقره دومین مورد خیلی داخلی 
%                \پایان{فقرات}
%            \فقره مورد دوم داخلی 
%        \پایان{فقرات}
%    \فقره مورد دوم 
%\پایان{فقرات}

% لیست موارد شماره‌گذاری شده:
%\شروع{شمارش}
%    \فقره یک مورد 
%    \فقره یک مورد دیگر 
%    \فقره باز هم یک مورد دیگر
%\پایان{شمارش}

% برای نوشتن یک فرمول چند خطی که یک شماره بخوره. برای حالت‌های دیگر به سایت زیر مراجعه کنید:
% https://www.overleaf.com/learn/latex/Aligning_equations_with_amsmath
%\begin{align} 
%    \begin{split}
%        (x+y)^3 &= (x+y)^2(x+y) \\
%                &= (x^2+2xy+y^2)(x+y) \\
%                &= (x^3+2x^2y+xy^2) + (x^2y+2xy^2+y^3) \\
%                &= x^3+3x^2y+3xy^2+y^3
%    \end{split}                                 
%\end{align}

% نوشتن ماتریس و آرایه دوبعدی
%\begin{align}
%    A = 
%    \begin{bmatrix}
%        A_{11} & A_{21} \\
%        A_{21} & A_{22}
%    \end{bmatrix}
%\end{align}

% اضافه کردن تصویر با شرح. برای اضافه کردن تصویر بدون شرح، فقط دستور درج‌تصویر را استفاده کنید.
%\شروع{شکل}[h!]
%\تنظیم‌ازوسط
%  \درج‌تصویر[پهنا=0.5\پهنای‌متن]{img.jpg}
%\شرح{شرح تصویر}
%\برچسب{تصویر:۱}
%\پایان{شکل}
%\پایان{زیرمسائل}

%کد زیر را ببینید.
%%\lstinputlisting{source_filename.py}
%\begin{latin}
%\begin{lstlisting}
%# Import the modules
%import sys
%import random
%
%ans = True
%
%while ans:
%question = raw_input("Ask the magic 8 ball a question: (press enter to quit) ")
%
%answers = random.randint(1,8)
%
%if question == "":
%sys.exit()
%
%elif answers == 1:
%print "It is certain"
%
%elif answers == 2:
%print "Outlook good"
%
%elif answers == 3:
%print "You may rely on it"
%
%elif answers == 4:
%print "Ask again later"
%
%elif answers == 5:
%print "Concentrate and ask again"
%
%elif answers == 6:
%print "Reply hazy, try again"
%
%elif answers == 7:
%print "My reply is no"
%
%elif answers == 8:
%print "My sources say no"
%\end{lstlisting}
%\end{latin}
